\section{Clinical rule interpretability}
\label{sec:rules}

Tables~\ref{tab:rules_heart_cleveland} and~\ref{tab:rules_thyroid_sick} show two
deployed MO-Fuzzy rule bases, one per dataset type: development (Cleveland) and
held-out (thyroid). To avoid choosing favourable examples, each is taken from the
outer fold whose test AUROC is closest to the median over the 25 folds of that dataset.
The model is read like a point score. Start from the baseline log-odds $w_0$ and add
$w_r$ for every rule whose conditions hold, scaled by its firing strength when a
condition holds only partly. Then convert the sum to a probability with the logistic
function. Fig.~\ref{fig:mf} shows the learnt membership functions in original units.

\begin{table}[t]
\centering
\scriptsize
\caption{Deployed MO-Fuzzy rule base for Cleveland (seed 0, fold 2: the fold whose test AUROC is closest to the median over 25 folds; test AUROC 0.902, Brier 0.131). Risk score: $\mathrm{logit}\,\hat p = w_0 + \sum_r w_r \beta_r(x)$. Membership-function centres in original units.}
\label{tab:rules_heart_cleveland}
\begin{adjustbox}{max width=\linewidth}
\begin{tabular}{p{0.6cm}p{5.2cm}r}
\toprule
Rule & IF (antecedents) & $w_r$ (log-odds) \\
\midrule
$w_0$ & baseline & -1.977 \\
R1 & chol is High AND thal=normal is No & +2.732 \\
R2 & slope=up is No & +2.181 \\
R3 & ca is Low & -2.764 \\
R4 & sex=male is Yes AND fbs=t is No & +1.564 \\
R5 & cp=asympt is Yes & +1.870 \\
\midrule
\multicolumn{3}{l}{Low / Medium / High centres:} \\
\multicolumn{3}{l}{\quad chol: 206, 225, 286} \\
\multicolumn{3}{l}{\quad ca: 0.0189, 0.169, 0.782} \\
\bottomrule
\end{tabular}
\end{adjustbox}
\end{table}

\begin{table}[t]
\centering
\scriptsize
\caption{Deployed MO-Fuzzy rule base for Thyroid (seed 3, fold 3: the fold whose test AUROC is closest to the median over 25 folds; test AUROC 0.976, Brier 0.023). Risk score: $\mathrm{logit}\,\hat p = w_0 + \sum_r w_r \beta_r(x)$. Membership-function centres in original units.}
\label{tab:rules_thyroid_sick}
\begin{adjustbox}{max width=\linewidth}
\begin{tabular}{p{0.6cm}p{5.2cm}r}
\toprule
Rule & IF (antecedents) & $w_r$ (log-odds) \\
\midrule
$w_0$ & baseline & -5.381 \\
R1 & referral\_source=SVI is No & -1.689 \\
R2 & T3 is Low & +6.472 \\
R3 & TT4 is High & +3.855 \\
\midrule
\multicolumn{3}{l}{Low / Medium / High centres:} \\
\multicolumn{3}{l}{\quad T3: 1.07, 1.45, 2.31} \\
\multicolumn{3}{l}{\quad TT4: 90.8, 151, 161} \\
\bottomrule
\end{tabular}
\end{adjustbox}
\end{table}


The Cleveland model has \resRulesNCleveland{} rules. Asymptomatic chest pain, a
non-upsloping ST segment, male sex without raised fasting blood sugar, and high serum
cholesterol without a normal thallium scan each add risk. A low number of major vessels
coloured by fluoroscopy (\texttt{ca}) lowers it. All are standard variables of the
study that produced the data~\cite{detrano1989}, recorded routinely, so the score can be
computed at the bedside without software. Whether the rules are clinically acceptable
must still be judged by domain experts.

The thyroid model has \resRulesNThyroid{} rules, and a single rule on a low T3
concentration carries most of the weight. One rule refers to the referral source rather
than to the patient's physiology: patients not referred through the \texttt{SVI} source
have lower risk. A rule like this records how patients reached the laboratory, not why
they were ill. It is a proxy of the kind discussed under leakage~\cite{kaufman2012}
and may not transport to other settings. Black-box models can exploit such a variable just as
well, but there it stays hidden. A model with \resRulesNThyroid{} rules makes the dependence visible
at a glance and lets a domain expert remove the variable and refit. This is the
practical meaning of interpretability that our literal count is meant to
approximate~\cite{rudin2019}.

\begin{figure}[t]
\centering
\includegraphics[width=0.8\linewidth]{mf_heart_cleveland}\\
\includegraphics[width=0.8\linewidth]{mf_thyroid_sick}
\caption{Learnt Low/Medium/High membership functions (original units) of the continuous
variables used by the rule bases in Tables~\ref{tab:rules_heart_cleveland} (top) and
\ref{tab:rules_thyroid_sick} (bottom). The partitions are strong (memberships sum to
one), so each value has a unique linguistic reading.}
\label{fig:mf}
\end{figure}
